% GENERATED FILE. Edit canonical lessons, then run npm run generate:books.
% canonical-source-hash: fnv1a64:25f0642b8ad0ad0e
% canonical-lessons: JA-C520-kigakari, JA-C520-kokorobosoi, JA-C520-omoigakenai, JA-C520-yosougai, JA-C520-kakujitsu

\chapter{Worries and Surprises}
\label{ch:ja-worries-and-surprises}
\hlchaptermodality{520}

\begin{chapteropening}
\textbf{By the end of this chapter:} \emph{I can say what to do if something unexpected happens.}

Complete the last lesson of chapter 520: I can say what to do if something unexpected happens.
\end{chapteropening}

\section[\texorpdfstring{\ja{きがかり} (kigakari)}{きがかり (kigakari)}]{\ja{きがかり} (kigakari) — a worry; something on one's mind}

\label{lesson:JA-C520-kigakari}



\begin{quote}
Before the new one: say the Japanese for to betray, then the Japanese for safe (as a choice).
\end{quote}

\subsection*{You'll want to know: \ja{きがかり}}
\textbf{\ja{きがかり}} — \emph{kigakari} — \textquotedblleft{}a worry; something on one's mind\textquotedblright{}.

\begin{quote}
\textbf{\ja{もし} \ja{あめが} \ja{ふったら、ははは} \ja{こどもの} \ja{ことが} \ja{きがかりでしょう。}} — \emph{moshi ame ga futtara, haha wa kodomo no koto ga kigakari deshō} — \textquotedblleft{}If it rains, my mother will be worried about the children.\textquotedblright{}
\end{quote}

A na-adjective: \ja{な} before a noun.

\begin{cousinweb}
\textbf{\ja{きがかり}} (\emph{kigakari}) mixes the two layers, native and Sino-Japanese: Sino-Japanese \emph{ki}, \textquotedblleft{}spirit\textquotedblright{} (Middle Chinese \emph{khj+jH}), borrowed from Chinese, and native \emph{gakari}, \textquotedblleft{}hanging\textquotedblright{}, the voiced form of \emph{kakari}, from \emph{kakaru}, \textquotedblleft{}to hang\textquotedblright{} (JMdict, KANJIDIC2; Middle Chinese in Baxter's transcription).
\end{cousinweb}

\begin{grammarlens}[title={What you've built}]
One more word for saying what happens if.
\end{grammarlens}

\subsection*{Your turn}
\begin{itemize}
\raggedright
  \item \emph{Say it:} \emph{kigakari}
  \item \emph{Say it:} \emph{kigakari}, once more
  \item \emph{Say it:} \emph{moshi ame ga futtara, haha wa kodomo no koto ga kigakari deshō} — If it rains, my mother will be worried about the children.
  \item \emph{Recall:} say the Japanese for to betray, then the Japanese for safe (as a choice), then say \emph{kigakari} again
\end{itemize}

\begin{tcolorbox}[breakable,colback=teal!4,colframe=teal!35!black,title={Before you move on}]
What does \ja{きがかり} mean? (\textquotedblleft{}A worry; something on one's mind\textquotedblright{}.) Say it once more.
\end{tcolorbox}
\section[\texorpdfstring{\ja{こころぼそい} (kokorobosoi)}{こころぼそい (kokorobosoi)}]{\ja{こころぼそい} (kokorobosoi) — helpless; lonely and uneasy}

\label{lesson:JA-C520-kokorobosoi}



\begin{quote}
Before the new one: say the Japanese for safe (as a choice), then the Japanese for a worry.
\end{quote}

\subsection*{You'll want to know: \ja{こころぼそい}}
\textbf{\ja{こころぼそい}} — \emph{kokorobosoi} — \textquotedblleft{}helpless; lonely and uneasy\textquotedblright{}.

\begin{quote}
\textbf{\ja{もし} \ja{ひとりで} \ja{いったら、いもうとは} \ja{こころぼそいでしょう。}} — \emph{moshi hitori de ittara, imōto wa kokorobosoi deshō} — \textquotedblleft{}If my younger sister went on her own, she would feel lonely and uneasy.\textquotedblright{}
\end{quote}

An i-adjective.

\begin{cousinweb}
\textbf{\ja{こころぼそい}} (\emph{kokorobosoi}) is a native Japanese compound (wago), inherited rather than borrowed from Chinese: built from \emph{kokoro}, \textquotedblleft{}heart\textquotedblright{} and \emph{bosoi}, the voiced form of \emph{hosoi}, \textquotedblleft{}thin\textquotedblright{} (JMdict, KANJIDIC2).
\end{cousinweb}

\begin{grammarlens}[title={What you've built}]
One more word for saying what happens if.
\end{grammarlens}

\subsection*{Your turn}
\begin{itemize}
\raggedright
  \item \emph{Say it:} \emph{kokorobosoi}
  \item \emph{Say it:} \emph{kokorobosoi}, once more
  \item \emph{Say it:} \emph{moshi hitori de ittara, imōto wa kokorobosoi deshō} — If my younger sister went on her own, she would feel lonely and uneasy.
  \item \emph{Recall:} say the Japanese for safe (as a choice), then the Japanese for a worry, then say \emph{kokorobosoi} again
\end{itemize}

\begin{tcolorbox}[breakable,colback=teal!4,colframe=teal!35!black,title={Before you move on}]
What does \ja{こころぼそい} mean? (\textquotedblleft{}Helpless; lonely and uneasy\textquotedblright{}.) Say it once more.
\end{tcolorbox}
\section[\texorpdfstring{\ja{おもいがけない} (omoigakenai)}{おもいがけない (omoigakenai)}]{\ja{おもいがけない} (omoigakenai) — unexpected}

\label{lesson:JA-C520-omoigakenai}



\begin{quote}
Before the new one: say the Japanese for a worry, then the Japanese for helpless.
\end{quote}

\subsection*{You'll want to know: \ja{おもいがけない}}
\textbf{\ja{おもいがけない}} — \emph{omoigakenai} — \textquotedblleft{}unexpected\textquotedblright{}.

\begin{quote}
\textbf{\ja{もし} \ja{おもいがけない} \ja{ことが} \ja{あったら、しらせて} \ja{ください。}} — \emph{moshi omoigakenai koto ga attara, shirasete kudasai} — \textquotedblleft{}If something unexpected happens, let me know.\textquotedblright{}
\end{quote}

An i-adjective.

\begin{cousinweb}
\textbf{\ja{おもいがけない}} (\emph{omoigakenai}) is a native Japanese word (wago), inherited rather than borrowed from Chinese; its kanji was attached later as a kun'yomi, and its older origin is not settled (JMdict, KANJIDIC2).
\end{cousinweb}

\begin{grammarlens}[title={What you've built}]
One more word for saying what happens if.
\end{grammarlens}

\subsection*{Your turn}
\begin{itemize}
\raggedright
  \item \emph{Say it:} \emph{omoigakenai}
  \item \emph{Say it:} \emph{omoigakenai}, once more
  \item \emph{Say it:} \emph{moshi omoigakenai koto ga attara, shirasete kudasai} — If something unexpected happens, let me know.
  \item \emph{Recall:} say the Japanese for a worry, then the Japanese for helpless, then say \emph{omoigakenai} again
\end{itemize}

\begin{tcolorbox}[breakable,colback=teal!4,colframe=teal!35!black,title={Before you move on}]
What does \ja{おもいがけない} mean? (\textquotedblleft{}Unexpected\textquotedblright{}.) Say it once more.
\end{tcolorbox}
\section[\texorpdfstring{\ja{よそうがい} (yosōgai)}{よそうがい (yosōgai)}]{\ja{よそうがい} (yosōgai) — beyond expectation; unforeseen}

\label{lesson:JA-C520-yosougai}



\begin{quote}
Before the new one: say the Japanese for helpless, then the Japanese for unexpected.
\end{quote}

\subsection*{You'll want to know: \ja{よそうがい}}
\textbf{\ja{よそうがい}} — \emph{yosōgai} — \textquotedblleft{}beyond expectation; unforeseen\textquotedblright{}.

\begin{quote}
\textbf{\ja{もし} \ja{よそうがいの} \ja{あめが} \ja{ふったら、しあいは} \ja{ちゅうしでしょう。}} — \emph{moshi yosōgai no ame ga futtara, shiai wa chūshi deshō} — \textquotedblleft{}If unexpected rain falls, the match will probably be stopped.\textquotedblright{}
\end{quote}

A na-adjective: \ja{な} before a noun.

\begin{cousinweb}
\textbf{\ja{よそうがい}} (\emph{yosōgai}) is Sino-Japanese (kango): three Chinese characters read with their on'yomi, readings borrowed from Chinese. They are \emph{yo}, \textquotedblleft{}beforehand\textquotedblright{} (Middle Chinese \emph{yoH}), \emph{sō}, \textquotedblleft{}concept\textquotedblright{} (\emph{sjangX}) and \emph{gai}, \textquotedblleft{}outside\textquotedblright{} (\emph{ngwajH}) (KANJIDIC2; Middle Chinese in Baxter's transcription).
\end{cousinweb}

\begin{grammarlens}[title={What you've built}]
One more word for saying what happens if.
\end{grammarlens}

\subsection*{Your turn}
\begin{itemize}
\raggedright
  \item \emph{Say it:} \emph{yosōgai}
  \item \emph{Say it:} \emph{yosōgai}, once more
  \item \emph{Say it:} \emph{moshi yosōgai no ame ga futtara, shiai wa chūshi deshō} — If unexpected rain falls, the match will probably be stopped.
  \item \emph{Recall:} say the Japanese for helpless, then the Japanese for unexpected, then say \emph{yosōgai} again
\end{itemize}

\begin{tcolorbox}[breakable,colback=teal!4,colframe=teal!35!black,title={Before you move on}]
What does \ja{よそうがい} mean? (\textquotedblleft{}Beyond expectation; unforeseen\textquotedblright{}.) Say it once more.
\end{tcolorbox}
\section[\texorpdfstring{\ja{かくじつ} (kakujitsu)}{かくじつ (kakujitsu)}]{\ja{かくじつ} (kakujitsu) — sure; reliable}

\label{lesson:JA-C520-kakujitsu}



\begin{quote}
Before the new one: say the Japanese for unexpected, then the Japanese for beyond expectation.
\end{quote}

\subsection*{You'll want to know: \ja{かくじつ}}
\textbf{\ja{かくじつ}} — \emph{kakujitsu} — \textquotedblleft{}sure; reliable\textquotedblright{}.

\begin{quote}
\textbf{\ja{もし} \ja{でんしゃで} \ja{いったら、くるまより} \ja{かくじつです。}} — \emph{moshi densha de ittara, kuruma yori kakujitsu desu} — \textquotedblleft{}If you go by train, it is more reliable than by car.\textquotedblright{}
\end{quote}

A na-adjective: \ja{な} before a noun.

\begin{cousinweb}
\textbf{\ja{かくじつ}} (\emph{kakujitsu}) is Sino-Japanese (kango): two Chinese characters read with their on'yomi, readings borrowed from Chinese. They are \emph{kaku}, \textquotedblleft{}assurance\textquotedblright{} (Middle Chinese \emph{khaewk}) and \emph{jitsu}, \textquotedblleft{}reality\textquotedblright{} (\emph{zyit}) (KANJIDIC2; Middle Chinese in Baxter's transcription).
\end{cousinweb}

\begin{grammarlens}[title={What you've built}]
One more word for saying what happens if.
\end{grammarlens}

\subsection*{Your turn}
\begin{itemize}
\raggedright
  \item \emph{Say it:} \emph{kakujitsu}
  \item \emph{Say it:} \emph{kakujitsu}, once more
  \item \emph{Say it:} \emph{moshi densha de ittara, kuruma yori kakujitsu desu} — If you go by train, it is more reliable than by car.
  \item \emph{Recall:} say the Japanese for unexpected, then the Japanese for beyond expectation, then say \emph{kakujitsu} again
\end{itemize}

\begin{tcolorbox}[breakable,colback=teal!4,colframe=teal!35!black,title={Before you move on}]
What does \ja{かくじつ} mean? (\textquotedblleft{}Sure; reliable\textquotedblright{}.) Say it once more.
\end{tcolorbox}
